\section*{अध्याय \ref{ch:b2:polymer-synthesis} --- बहुलक: संश्लेषण, संरचना और गुणधर्म}

\begin{solution}{exo:b2:polymer-synthesis:1}\emergencystretch=3em
पॉलिएथीन \ce{-[CH2-CH2]-}; पॉलिप्रोपीन \ce{-[CH2-CH(CH3)]-}; पॉली(वाइनिल क्लोराइड)
\ce{-[CH2-CHCl]-}; पॉलिस्टाइरीन \ce{-[CH2-CH(C6H5)]-}; नायलॉन-6,6 \ce{-[NH(CH2)6NH-CO(CH2)4CO]-}।
\end{solution}

\begin{solution}{exo:b2:polymer-synthesis:2}
PET: चरण-वृद्धि (पॉलिएस्टर, जल मुक्त होता है)। पॉलिस्टाइरीन: शृंखला-वृद्धि। नायलॉन-6,6: चरण-वृद्धि।
पॉली(मेथिल मेथैक्रिलेट): शृंखला-वृद्धि (एक \ce{C=C} \omterm{def:b2:polymer-synthesis:polymer}{एकलक})। हाइड्रॉक्सी अम्ल से पॉली(लैक्टिक अम्ल):
चरण-वृद्धि, एक \ce{A-B} \omterm{def:b2:polymer-synthesis:polymer}{एकलक}।
\end{solution}

\begin{solution}{exo:b2:polymer-synthesis:3}
$\bar M_n = 1000 \times \qty{104.15}{g/mol} \approx \qty{104}{kg/mol}$ (स्टाइरीन \ce{C8H8},
\qty{104.152}{g/mol})।
% ledger: aw:C, aw:H
\end{solution}

\begin{solution}{exo:b2:polymer-synthesis:4}
\omterm{def:b2:polymer-synthesis:tacticity}{सिंडियोटैक्टिक}। हाँ: नियमित शृंखला क्रिस्टलकों में संकुलित हो सकती है।
\end{solution}

\begin{solution}{exo:b2:polymer-synthesis:5}
द्रव्यमान अंश 0.5 और 0.5। $1/\bar M_n = 0.5/10 + 0.5/100 = 0.055$, $\bar M_n \approx
\qty{18.2}{kg/mol}$; $\bar M_w = 0.5 \times 10 + 0.5 \times 100 = \qty{55}{kg/mol}$; $\text{\DH} \approx
3.0$।
\end{solution}

\begin{solution}{exo:b2:polymer-synthesis:6}
$1/(1-0.98) = 50$; $1/(1-0.995) = 200$।
\end{solution}

\begin{solution}{exo:b2:polymer-synthesis:7}
$r = 1/1.01$; $p = 1$ पर $\bar X_n = (1+r)/(1-r) = (1.01 + 1)/(1.01 - 1) = 201$।
\end{solution}

\begin{solution}{exo:b2:polymer-synthesis:8}
$R_p \propto [\mathrm I]^{1/2}$: $\sqrt2 \approx 1.41$ से गुणा। $\nu \propto [\mathrm I]^{-1/2}$:
$\sqrt 2$ से भाग, अपने पूर्व मान का लगभग 0.71 गुना।
\end{solution}

\begin{solution}{exo:b2:polymer-synthesis:9}
शुद्ध PVC का काँच संक्रमण कमरे के ताप से ऊपर है: एक दृढ़ काँच। छोटे सुघट्यकारी
अणु शृंखलाओं के बीच बैठते हैं, उन्हें अलग करते हैं और खंडों को निम्नतर ताप पर गति करने देते हैं: $T_g$
कमरे के ताप से नीचे गिर जाता है, और उपयोग में पदार्थ रबर जैसा और लचीला होता है।
\end{solution}

\begin{solution}{exo:b2:polymer-synthesis:10}
$\dfrac{\mathrm d}{\mathrm dx}\bigl(x p^{\,x-1}\bigr) = p^{\,x-1}(1 + x\ln p)$, जो $x = -1/\ln p$ के लिए शून्य है
(एक उच्चिष्ठ: चिह्न $+$ से $-$ हो जाता है)। $p$ के 1 के निकट होने पर $\ln p \approx -(1-p)$, अतः $x \approx
1/(1-p) = \bar X_n$।
\end{solution}

\begin{solution}{exo:b2:polymer-synthesis:11}
चरण 1, \omterm{def:b2:acyl-substitution:transesterification}{पार-एस्टरीकरण}: \ce{C6H4(COOCH3)2 + 2HOCH2CH2OH -> C6H4(COOCH2CH2OH)2 + 2CH3OH}, मेथेनॉल
आसवित करके हटाया जाता है। चरण 2: हर डाइएस्टर अणु के दो \ce{CH2CH2OH} सिरे हैं; एक सिरा किसी दूसरे अणु के
एस्टर पर आक्रमण करता है और एथेन-1,2-डाइऑल का एक अणु मुक्त करता है, जिसे पंप करके हटा दिया जाता है। \omterm{def:b2:polymer-synthesis:polymer}{एकलक}
दोनों प्रकार के समूह सही अनुपात में धारण करता है, जैसे कोई \ce{A-B} \omterm{def:b2:polymer-synthesis:polymer}{एकलक} करता है: संतुलन $r = 1$
उसमें अंतर्निहित है, और डाइऑल को हटाना $p$ को 1 की ओर ले जाता है।
\end{solution}

\begin{solution}{exo:b2:polymer-synthesis:12}
समान-मोलर भरण: अनुपात $= (2.0 + 1)/(1 + 0.5) = 2$: \omterm{def:b2:polymer-synthesis:polymer}{एकलक} 1 की इकाइयाँ दुगुनी। $r_1 = r_2 =
0$ के साथ अनुपात $= [\mathrm M_1][\mathrm M_2]/([\mathrm M_2][\mathrm M_1]) = 1$, भरण चाहे जो हो: हर
मूलक केवल दूसरा \omterm{def:b2:polymer-synthesis:polymer}{एकलक} जोड़ता है, अर्थात् एकांतरी \omterm{def:b2:polymer-synthesis:copolymer}{सहबहुलक}।
\end{solution}

\begin{solution}{pb:b2:polymer-synthesis:1}
\textbf{1.} \ce{H2N(CH2)6NH2} और \ce{HOOC(CH2)4COOH}।
\textbf{2.} प्रति कड़ी एक जल:
\[\mathrm{{-}COOH + H_2N{-} \longrightarrow {-}CO{-}NH{-} + H_2O}.\]
\textbf{3.} लवण में ठीक प्रति डाइअम्ल एक डाइऐमीन होता है; दो अलग द्रवों या ठोसों को तौलने से
कैरोदर्स समीकरण द्वारा माँगी गई यथातथता कभी नहीं मिलती।
\textbf{4.} \ce{-[NH(CH2)6NH-CO(CH2)4CO]-}, \ce{C12H22N2O2}, \qty{226.32}{g/mol}।
\textbf{5.} \omterm{def:b2:polymer-synthesis:polymer}{पुनरावृत्त इकाई} में दो एकलक-व्युत्पन्न इकाइयाँ हैं: $M_0 = 226.32/2 = \qty{113.16}{g/mol}$।
\textbf{6.} डाइऐमीन (6) और डाइअम्ल (6) के कार्बन।
\textbf{7.} $\bar X_n = 50$, $\bar M_n = 50 \times 113.16 \approx \qty{5.66}{kg/mol}$।
\textbf{8.} $\bar X_n = 100$, $\bar M_n \approx \qty{11.3}{kg/mol}$।
\textbf{9.} 99~\% समूहों के अभिक्रिया कर लेने पर औसत शृंखला में केवल 100 इकाइयाँ होती हैं, जो
मज़बूत रेशे के लिए बहुत छोटी है; पूर्ण \omterm{def:b2:continuous-reactors:conversion}{रूपांतरण} की ओर हर अगला कदम शृंखला की लंबाई को दुगुना या तिगुना कर देता है।
\textbf{10.} $\bar X_n = 10\,000/113.16 = 88.4$; $p = 1 - 1/88.4 \approx 0.989$।
\textbf{11.} जल हटाकर: गलित को उसके गलनांक से काफ़ी ऊपर गर्म किया जाता है, अंत में घटे हुए
दाब पर, ताकि साम्य ऐमाइड की ओर खिसकता रहे।
\textbf{12.} कोई एकक्रियाशील अशुद्धि शृंखला-सिरों को बंद कर देती है, और कोई भी अशुद्धि 1\,:\,1 संतुलन को बिगाड़ देती है।
\textbf{13.} $r = 1/1.01 \approx 0.990$।
\textbf{14.} $\bar X_n = (1+r)/(1-r) = 201$; $\bar M_n \approx 201 \times 113.16 \approx
\qty{22.7}{kg/mol}$।
\textbf{15.} $\bar X_n = 1.9901/(1.9901 - 2 \times 0.9901 \times 0.99) \approx 1.9901/0.0297 \approx
67$।
\textbf{16.} ऐमीन समूह: जब अम्ल समूह समाप्त हो जाते हैं, तो हर शृंखला \ce{NH2} पर समाप्त होती है।
\textbf{17.} यह किसी ऐमीन सिरे को ऐसे ऐमाइड में बदल देता है जो आगे अभिक्रिया नहीं कर सकता: यह शृंखलाओं को बंद करता है,
$\bar M_n$ को सीमित करता है और \omterm{def:b2:polymer-synthesis:polymer}{बहुलक} को फिर से पिघलाने पर आगे की वृद्धि रोकता है।
\textbf{18.} कताई और ढलाई के लिए पुनरुत्पादनीय गलित श्यानता बनाए रखने के लिए, जो अन्यथा
गलित में शृंखलाओं के अभिक्रिया करते रहने से बदलती रहती।
\textbf{19.} $\text{\DH} = 1 + p = 1.99$।
\textbf{20.} $\bar M_w = 1.99 \times 11.3 \approx \qty{22.5}{kg/mol}$।
\textbf{21.} संख्या अंश $n_1 = 1 - p = 0.01$, अणुओं का 1~\%; द्रव्यमान अंश $w_1 = (1-p)^2 =
10^{-4}$, द्रव्यमान का 0.01~\%।
\textbf{22.} $x = -1/\ln 0.99 \approx 99.5$ के निकट, अर्थात् $\bar X_n = 100$ के निकट (अभ्यास 10)।
\textbf{23.} $1000/226.32 = \qty{4.42}{mol}$ पुनरावृत्त इकाइयाँ, हर एक में दो ऐमाइड कड़ियाँ: जल के \qty{8.84}{mol},
$8.84 \times 18.015 \approx \qty{159}{g}$।
\textbf{24.} बहुत छोटी शृंखलाएँ दुर्बल रेशा देती हैं; बहुत लंबी शृंखलाएँ इतना श्यान गलित देती हैं कि उसे
स्पिनरेट के महीन छिद्रों से धकेला नहीं जा सकता।
\textbf{25.} $\bar X_n = 20\,000/113.16 = 176.7$ और $p = 1 - 1/176.7 = \boldsymbol{\approx 0.994}$।
\end{solution}
